\documentclass[pdfbookmarks,a4paper,11pt]{article}
\usepackage{dsekcommon}
\usepackage[utf8]{inputenc}
\usepackage{enumerate}
\usepackage{lastpage}
\usepackage{longtable}
\ifpdf
  \usepackage{thumbpdf}
\fi
\usepackage{calc}

\pagestyle{fancy}

\lhead{}
\chead{}
\rhead{}
\lfoot{}
\cfoot{\thepage\ (\nohyperpageref{LastPage})}
\rfoot{}

% Puts a full stop after the section number
\makeatletter
\def\@seccntformat#1{\csname the#1\endcsname.\quad}
\makeatother

\newcommand{\funktionar}[1]{%
  \subsection*{#1}\par
  \addcontentsline{toc}{subsection}{#1}}

\newlength{\itemcollength}
\setlength{\itemcollength}{40mm}
\newenvironment{reglemlista}{%
  \begin{list}{}{%
      \setlength{\labelwidth}{\itemcollength}%
      \setlength{\leftmargin}{\labelwidth + \labelsep}%
      \renewcommand{\makelabel}[1]{%
        \raisebox{0pt}[1ex][0pt]{%
          \makebox[\labelwidth][l]{%
            \parbox[t]{\itemcollength}{%
              \raggedright\hspace{0pt}##1}}}\hfill}%
      }}{%
  \end{list}}

%%%%%%%%%%%%% HEADER %%%%%%%%%%%%%

\title{Reglemente för D-sektionen inom TLTH\\ Organisationsnummer: 845003-2878}
\date{\today}

\setheader{REGLEMENTE}{VTM 2023}{\today}

%%%%%%%%%%%%% HEADER %%%%%%%%%%%%%
\begin{document}

\maketitle

%%%%%%%%%%%%% SECTION %%%%%%%%%%%%%
\section{Teknologia}

\begin{reglemlista}

	\item[Råsa]
	En lämplig tolkning av råsa är 0xF280A1. Detta kan tolkas som Pantone
	189~U eller 35,3\,\% svart.

	\item[Rosa Pantern]
	Med Rosa Pantern åsyftas den figur som skapades av Friz Freleng och David
	DePatie till Blake Edwards film ''The Pink Panther'' år 1963.

	\item[Råsa Digitalis]
	Med Råsa Digitalis menas en råsa fingerborgsblomma, på latin Digitalis Purpurea.

\end{reglemlista}


%%%%%%%%%%%%% SECTION %%%%%%%%%%%%%
\section{Hedersomnämnande}

\begin{reglemlista}

	\item[Hedersmedlemmar]
	Hedersmedlemmar är:
	\begin{itemize}
		\item Nina Reistad\\
		      \emph{Inlagd i reglementet på VTM1 2000. Inröstad 1993}
		\item Per-Henrik Rasmussen \\ \emph{inröstad 1994}
		\item Rune Kullberg\\
		      \emph{Invald på HTM1 2002 för sitt arbete för studentinflytande, hans handlingskraft i studentfrågor, och hans stora initiativförmåga som ordförande i Utbildningsnämnden för D}
		\item Nora Ekdahl\\
		      \emph{Invald VTM 2010 för sina insatser som studievägledare för C- och D-programmet.}
		\item Mats Cedervall \\
		      \emph{Invald på HTM1 2018 för sitt stora Intresse i sektionens lokaler och för att han som Husprefekt alltid har verkat för studenternas ökade inflytande i husstyrelsen}
		\item Jan Eric Larsson \\
			  \emph{Invald på HTM1 2022 för sina insatser och sitt stora engagemang som Inspektor under en lång tid.}
		\item Anu Uus \\
			  \emph{Invald på HTM1 2022 för sina insatser och sitt stora engagemang som Inspektrix under en lång tid.}
	\end{itemize}

	\item[Grundarna av D-sektionen inom TLTH]
	Grundarna av D-sektionen inom TLTH är:

	\begin{itemize}
		\item Lars Svensson
		\item Ulf Bengtsson
		\item Henrik Cronström
		\item Kristin Andersson
		\item Stefan Molund
		\item Jan Eric Larsson
	\end{itemize}

\end{reglemlista}


%%%%%%%%%%%%% SECTION %%%%%%%%%%%%%
\section{Strategiska mål och verksamhetsplan}

\begin{reglemlista}

	\item[Strategiska mål och verksamhetsplan]
	De strategiska målen ska innehålla strategiska mål som sektionen ska arbeta med
över längre tid. Verksamhetsplanen ska utgå från de strategiska målen och definiera konkreta mål som ska arbetas med under verksamhetsåret. Ett strategiskt mål
kan sträcka sig över flera år och ska varje år ha en konkret motsvarighet i verksamhetsplanen.

\end{reglemlista}

%%%%%%%%%%%%% SECTION %%%%%%%%%%%%%
\section{Policyer}

\begin{reglemlista}
	\item[Policyer]
    Sektionen har ett antal policyer som fungerar som förlängningar av reglementet gällande en särskild fråga. Notera att brott mot policyer kan medföra påföljder enligt TLTH:s “Policy för entledigande och avstängning”. En förteckning med sektionens policyer återfinns i policy för styrdokument.
\end{reglemlista}


%%%%%%%%%%%%% SECTION %%%%%%%%%%%%%
\section{Sektionsmöte}

\begin{reglemlista}

	\item[Utlysande]
	Utöver skriftligt anslag i sektionslokalen skall följande tillställas kallelse:
	\begin{itemize}
		\item Valberedningen
		\item Revisorerna
		\item Kårordföranden
		\item Inspektorn
	\end{itemize}
	
	\item[Förtydligande]
	Efter varje post står en siffra inom parentes, det är antal personer som kan sitta på en post, om det inte står någon siffra så är det erforderligt antal.

	\item[Vårterminsmöte]
	Under vårterminsmötet skall följande poster tillsättas:
	\begin{itemize}

		\item Teknikfokusansvarig (1)
		\item Jubileumsgeneral (1) (väljes året innan jubileet äger rum)
		\item Valberedningens ordförande (1)
		\item  Valberedningsrepresentanter som förväntas vara ordinarie medlem vid TLTH ett år framöver (Det totala antalet valberedningsrepresentanter får vid inget tillfälle överstiga 20.)
		\item Tackmästare (1)
		\item Studerandeskyddsombud (1)
	\end{itemize}

	\item[Höstterminsmöte ett]
	Under höstterminsmöte ett skall följande poster tillsättas:
	\begin{itemize}
		\item Valberedningsrepresentant från den lägsta årskursen (Det totala
antalet valberedningsrepresentanter får vid inget tillfälle överstiga
20.)
		\item Øverphøs (1)
		\item Trivselmästare (1)
	\end{itemize}

	\item[Valmöte]
	Under valmötet skall följande poster tillsättas:
	\begin{itemize}
		\item Ordförande (1)
		\item Vice Ordförande (1)
		\item Informationsansvarig (1)
		\item Skattmästare (1)
		\item Cafémästare (1)
		\item Sexmästare (1)
		\item Studierådsordförande (1)
		\item Näringslivsansvarig (1)
		\item Källarmästare (1)
		\item Aktivitetsansvarig (1)
        \item Processmästare (1)
		\item Vice Informationsansvarig (1)
		\item Vice Skattmästare (1)
		\item Vice Cafémästare (1)
		\item Vice Sexmästare (1)
		\item Vice Studierådsordförande (1)
		\item Vice Näringslivsansvarig (1)
		\item Vice Källarmästare (1)
		\item Vice Aktivitetsansvarig (1)
        \item Vice Processmästare (1)
		\item Revisor (2)
		\item Inspektor (1)
		\item Framtidsordförande (1)
		\item Valnämndsrepresentant TLTH (1)
		\item Alumniansvarig (1)
		\item Utedischoansvarig (1-2)
		\item Övermarskalk (1)
		\item Talman (1)
		\item Skattförman (1-4)
		\item Tackmästare (1)
	\end{itemize}

\end{reglemlista}

%%%%%%%%%%%%% SECTION %%%%%%%%%%%%%
\section{Valberedningen}

\begin{reglemlista}

	\item[Åligganden]
	Det åligger valberedningen
	\begin{attlista}
		\item valbereda kandidater till de poster som väljs under sektionsmöte med undantag för poster inom Valberedningen
		\item kontakta och informera samtliga kandidater om resultatet från valprocessen innan det
		väljande mötet
		\item skicka in förslag på vilka kandidater Valberedningen anser vara lämpliga i form av en
		handling till det väljande mötet
	\item skicka minst en valberedningsrepresentant till varje utskottsvalberedning.
	\end{attlista}
\end{reglemlista}

%%%%%%%%%%%%% SECTION %%%%%%%%%%%%%
\section{Styrelsen}

\begin{reglemlista}

	\item[Åligganden]
	Det åligger styrelsen
	\begin{attlista}
		\item inför sektionen ansvara för sektionens verksamhet
		\item förbereda och genomföra sektionsmöten
		\item verkställa och övervaka genomförandet av sektionsmötets beslut
		\item ansvara för sektionens medel
		\item bereda inkomna förslag
		\item handha sektionens korrespondens
		\item utarbeta förslag till budget och resultatdispositioner
		\item utarbeta förslag till verksamhetsplan
		\item senast sex veckor efter verksamhetsårets slut till revisor\-erna överlämna verksamhetsberättelse, protokoll och övriga handlingar revisorerna önskar ta del av
		\item en gång per termin informera medaljelekomittén om de funktionärer som styreslen anser skall mottaga medalj
		\item uträtta val som inte uträttats av sektionsmötet
		\item anordna överlämning för nästkommande styrelse
	\end{attlista}

	\item[Sammanträde]
	Sektionsstyrelsen sammanträder på kallelse av ordförande eller vid dennes frånvaro av vice ordförande. Ständigt adjungerad är, förutom de av stadgan och reglementet fastställda, även E-sektionens
	ordförande, två styrelseledarmöter från D-chip, sektionens kårkontakter, Trivselmästaren, Framtidsordförande samt Processmästaren.

\end{reglemlista}

%%%%%%%%%%%%% SECTION %%%%%%%%%%%%%
\section{Utskott}

\begin{reglemlista}

	\item[Åligganden]
	Det åligger utskotten
	\begin{attlista}
		\item söka erforderliga tillstånd
		\item följa av sektionen uppsatta riktlinjer
	\end{attlista}

	\item[Utskottsordföranden]
	Det åligger utskottsordföranden
	\begin{attlista}
		\item ansvara för utskottets verksamhet
		\item senast fyra veckor efter verksamhetsårets slut till styrelsen inlämna verksamhetsberättelse
		\item leda och fördela arbetet inom utskottet
		\item budgetera, redovisa och följa upp utskottets verksamhet
		\item upprätta bokslut
		\item sammankalla utskottet minst två gånger per verksamhetsår
		\item hålla styrelsen informerad om utskottets verksamhet
		\item utse erforderligt antal funktionärer utöver de av sektionen valda
		\item anordna överlämning för efterträdare
	\end{attlista}

	\item[Titulering]
	Alternativ titulering för utskottsordförande är mästare.

\end{reglemlista}

%%%%%%%%%%%%% SECTION %%%%%%%%%%%%%
\section{Cafémästeriet}

\begin{reglemlista}

	\item[Åligganden]
	Det åligger cafémästeriet
	\begin{attlista}
		\item ansvara för caféts verksamhet
		\item verka för att ett så brett cafésortiment som möjligt finns tillgängligt för sektionens medlemmar
		\item ansvara för daglig städning av iDét
		\item då det är praktiskt lämpligt, assistera D-shopen med försäljning av märken
	\end{attlista}

	\item[Sammansättning]
	Cafémästeriet består av
	\begin{itemize}
		\item Cafémästare
		\item Vice Cafémästare
		\item Dagsansvarig
		\item Stekare
        \item Brunchmästare
        \item Inventarieansvarig
        \item Bakis
		\item Funktionärer
	\end{itemize}

\end{reglemlista}

%%%%%%%%%%%%% SECTION %%%%%%%%%%%%%
\section{Näringslivsutskottet}

\begin{reglemlista}

	\item[Åligganden]
	Det åligger näringslivsutskottet
	\begin{attlista}
		\item verka för att sektionens goda relationer med näringslivet behålls och förbättras
		\item verka för att information om datateknik- och infocomutbildningarna vid LTH sprids till företag och näringsliv
		\item underlätta kontakten mellan sektionsmedlemmar och verkligheten
		\item vid behov söka sponsorer till sektionen
		\item genom kontakt med gamla sektionsmedlemmar kunna ge en inblick i verksamheten i intressanta företag
		\item bistå Nollningsutskottet i sponsorsökandet inför nollningen
	\end{attlista}

	\item[Sammansättning]
	Näringslivsutskottet består av
	\begin{itemize}
		\item Näringslivsansvarig
		\item Vice Näringlivsansvarig
		\item Mentorsansvarig
		\item Alumnigruppen
		\item Projektgruppen för Teknikfokus
		\item Funktionärer
	\end{itemize}

	\item[\textbf{Alumnigruppen}]

	\item[Åligganden]
	Det åligger alumnigruppen
	\begin{attlista}
		\item  upprätthålla en förteckning över utexaminerade D- och C-teknologer
		\item upprätthålla kontakten med utexaminerade D- och C-teknologer
		\item underhålla examensmonumentet
	\end{attlista}

	\item[Sammansättning]
	Alumnigruppen består av
	\begin{itemize}
		\item Alumniansvarig
		\item Alumnigruppsmedlem (erfoderligt antal)
	\end{itemize}

\end{reglemlista}
\textbf{Projektgruppen för Teknikfokus}
\begin{reglemlista}
	\item[Åligganden]
	Det åligger Projektgruppen för Teknikfokus

	\begin{attlista}
		\item på våren anordna arbetsmarknadsmässan Teknikfokus
	\end{attlista}

	\item[Sammansättning]
	Projektgruppen för Teknikfokus består av
	\begin{itemize}
		\item Teknikfokusansvarig
		\item Medlem i Projektgruppen för Teknikfokus (erfoderligt antal)
	\end{itemize}

\end{reglemlista}


%%%%%%%%%%%%% SECTION %%%%%%%%%%%%%
\section{Källarmästeriet}

\begin{reglemlista}

	\item[Åligganden]
	Det åligger källarmästeriet
	\begin{attlista}
		\item underhålla sektionens lokaler och inventarier
		\item handha bokningar av iDét
		\item hålla minst en Snickerboa per termin
	\end{attlista}

	%%% Togs bort vtm 2018 "resning reglemente" snickerboa bör dock förklaras
	% \item[Snickerboa]
	%   Under snickerboa ges sektionsmedlemmarna möjlighet att försköna och
	%   underhålla sektionens lokaler och inventarier. Till snickerboa ränner
	%   alla sektionsmedlemmar.

	\item[Sammansättning]
	Källarmästeriet består av
	\begin{itemize}
		\item Källarmästaren
		\item Vice Källarmästare
		\item Ljud- och ljusansvarig
		\item Trädgårdsmästare
        \item Bilansvarig
		\item Funktionärer
	\end{itemize}

\end{reglemlista}

%%%%%%%%%%%%% SECTION %%%%%%%%%%%%%
\section{Medaljelelekommittén}

Medaljelelekommittén delar ut medaljer till sektionsfunktionärer. Medaljelelekommittén leds av Övermarskalken.

\begin{reglemlista}

	\item[Åligganden]
	Det åligger medaljelelekommittén
	\begin{attlista}
		\item vid högtidliga tillfällen, minst två gånger per år, tillsammans med (på skiftesgasquen avgående) sektionsordförande dela ut medaljer och andra utmärkelser till funktionärer som gjort sig förtjänta av dylikt
		\item underhålla en lista över personer och utmärkelser på sektionens hemsida
		\item minst en gång per termin bjuda in sektionsordförande till möte, så att denna kan informera om styrelsens åsikter om funktionärers prestationer
	\end{attlista}

	\item[Sammansättning]
	Medaljelelekommittén består av
	\begin{itemize}
		\item Övermarskalken
		\item Inspektor
		\item Medaljelelekommittémedlemmar
	\end{itemize}
\end{reglemlista}

%%%%%%%%%%%%% SECTION %%%%%%%%%%%%%
\section{Aktivitetsutskottet}

\begin{reglemlista}

	\item[Åligganden]
	Det åligger Aktivitetsutskottet
	\begin{attlista}
		\item anordna för sektionens medlemmar trevliga evenemang
		\item deltaga i sångarstriden
		\item verka för att sektionen har lag med i av andra studentorganisationer anordnade aktiviteter, såsom Tandemstafetten och DÖMD
		\item informera om större arrangemang på andra orter
	\end{attlista}

	\item[Sammansättning]
	Aktivitetsutskottet består av
	\begin{itemize}
		\item Aktivitetsansvarig
		\item Vice Aktivitetsansvarig
		\item Utedischoansvarig
		\item Idrottsförman
		\item Karnevalsansvarig
		\item LANparty-ansvarig
		\item Sångarstridsförmän
		\item Tandemgeneral
		\item Semesterfirare
		\item D-sportare
		\item Nöjesförman
		\item Coach
		\item Funktionärer
	\end{itemize}


%%%%%%%%%%%%% SECTION %%%%%%%%%%%%%
\section{Informationsutskottet}

\begin{reglemlista}

	\item[Åligganden]
	Det åligger Informationsutskottet
	\begin{attlista}
		\item utge sektionens medier
		\item ansvara för sektionens anslagstavlor
		\item sända kopior av sektionens publikationer till lämpliga instanser, t~ex övriga sektioner
		\item upprätthålla sektionens arkiv
	\end{attlista}

	\item[Sammansättning]
	Informationsutskottet består av
	\begin{itemize}
		\item Informationsansvarig
		\item Vice Informationsansvarig
		\item Arkivarie
		\item Artist
		\item Fotograf
		\item Filmare
    \item Redaktionen
		\item Webmaster
    \item Redaktör
		\item Funktionärer
    \item D-shopen
	\end{itemize}


    \item[\textbf{Redaktionen}]

    \item[Åliggande]
    Det åligger Redaktionen
    \begin{attlista}
	\item framställa publikationer för sektionen
    \end{attlista}

    \item[Sammansättning]
    Redaktionen består av 
    \begin{itemize}
        \item Redaktör
        \item Journalist
    \end{itemize}

	\item[\textbf{DWWW}]

	\item[Åligganden]
	Det åligger DWWW
	\begin{attlista}
		\item underhålla och utveckla sektionens WWW-sidor
	\end{attlista}

	\item[Sammansättning]
	DWWW består av
	\begin{itemize}
		\item DWWW-ansvarig
		\item DWWW-medlem
	\end{itemize}

\end{reglemlista}


	\item[\textbf{D-shopen}]

	\item[Åligganden]
	Det åligger D-shopen
	\begin{attlista}
		\item handha inköp och försäljning av olika sektionsprofilerade produkter samt dylikt.
		\item förnya sortimentet då så krävs.
	\end{attlista}
	\item[Sammansättning]
	D-shopen består av
	\begin{itemize}
		\item Shopaholic
		\item Märkvärdig
		\item Skald
	\end{itemize}
\end{reglemlista}


%%%%%%%%%%%%% SECTION %%%%%%%%%%%%%
\section{Sexmästeriet}

\begin{reglemlista}

	\item[Åligganden]
	Det åligger sexmästeriet
	\begin{attlista}
		\item  anordna sittningar, pubar samt eftersläpp.
	\end{attlista}

	\item[Pubmästeriet]
	Pubmästeriet ansvarar för sektionens pubverksamhet.

	\item[Sammansättning]
	Sexmästeriet består av
	\begin{itemize}
		\item Sexmästare
		\item Vice Sexmästare
		\item Köksmästare
		\item Vice köksmästare
		\item Pubmästare
		\item Vice Pubmästare
		\item Barmästare
		\item Vice Barmästare
		\item Sektionskock
		\item Sångförman
		\item Preferensmästare
		\item Ölförman
		\item Hovmästare
		\item Vinförman
		\item Funktionärer
        \item TappaD
	\end{itemize}
\end{reglemlista}

%%%%%%%%%%%%% SECTION %%%%%%%%%%%%%
\section{Skattmästeriet}

\begin{reglemlista}

	\item[Åligganden]
	Det åligger skattmästeriet
	\begin{attlista}
		\item handha sektionens bokföring
		\item handha sektionens transaktioner
		\item skriva budgetförslag
		\item samordna sektionens ekonomiska verksamhet
		\item följa och bevaka sektionens likviditet
	\end{attlista}

	\item[Sammansättning]
	Skattmästeriet består av
	\begin{itemize}
		\item Skattmästare
		\item Vice Skattmästare
		\item Skattförman
		\item Funktionärer
	\end{itemize}
\end{reglemlista}

%%%%%%%%%%%%% SECTION %%%%%%%%%%%%%
\section{Nollningsutskottet}

\begin{reglemlista}

	\item[Åligganden]
	Det åligger Nollningsutskottet
	\begin{attlista}
		\item med hjälp av sektionen planera och genomföra nollningen
		\item tillse att tredje part ej drabbas
		\item ansvara för att aktiviteter under nollningen visar positiv bild av D-sektionen och Teknologkåren
		\item vara goda representanter för D-sektionen och Teknologkåren under nollningen
		\item inom Staben utse ett ställföreträdande Øverphøs
	\end{attlista}

	\item[Sammansättning]
	Nollningsutskottet består av
	\begin{itemize}
		\item Øverphøs
		\item Stabsmedlemmar
		\item \O verpeppare
		\item Peppare
		\item Phaddrar
		\item Pluggphaddrar
        \item Funktionärer
	\end{itemize}

	\item[Staben]
	Staben består av
	\begin{itemize}
		\item Øverphøs
		\item Stabsmedlemmar
	\end{itemize}

\end{reglemlista}

%%%%%%%%%%%%% SECTION %%%%%%%%%%%%%
\section{Studierådet}

\begin{reglemlista}

	\item[Åligganden]
	Det åligger studierådet
	\begin{attlista}
		\item ansvara för sektionens studiebevakning
		\item välja
		\begin{itemize}
			\item Studierådssekreterare
			\item Studentrepresentant, Institutionsstyrelse
			\item Studentrepresentant, Programledning
			\item Årskursrepresentant, D1
			\item Årskursrepresentant, C1
			\item Årskursrepresentant, D2
			\item Årskursrepresentant, C2
			\item Årskursrepresentant, D3
			\item Årskursrepresentant, C3
			\item Årskursrepresentant, VR/AR1
			\item Årskursrepresentant, VR/AR2
            \item Årskursrepresentant, Specialisering
			\item Funktionärer
		\end{itemize}
		
		\item föreslå representanter till TLTH:s externa organ
	\end{attlista}

	\item[Sammanträde]
	Studierådet sammanträder på kallelse av Studierådsordföranden eller vid
	dennes frånvaro av Vice Studierådsordföranden. Rätt att begära utlysning av
	studierådsmöte har
	\begin{enumerate}
		\item minst 10 sektionsmedlemmar
		\item revisor
		\item Inspektor
	\end{enumerate}
	
		\item[Sammansättning]
	Studierådet består av
	\begin{itemize}
		\item Studierådsordförande
		\item Vice Studierådsordförande
		\item Studierådssekreterare
		\item Studentrepresentant, Institutionsstyrelse
		\item Studentrepresentant, Programledning
		\item Årskursrepresentant, D1
		\item Årskursrepresentant, C1
		\item Årskursrepresentant, D2
		\item Årskursrepresentant, C2
		\item Årskursrepresentant, D3
		\item Årskursrepresentant, C3
		\item Årskursrepresentant, VR/AR1
		\item Årskursrepresentant, VR/AR2
        \item Årskursrepresentant, Specialisering
		\item Funktionärer
	\end{itemize}

	\item[Kallelse]
	Kallelse till studierådsmöte samt föredragningslista skall senast fem
	dagar före sammanträdet anslås skriftligen i sektionslokalen där alla sektionsmedlemmar kan se.

\end{reglemlista}

\section{Trivselrådet}

\begin{reglemlista}

	\item[Åligganden]
	Det åligger Trivselrådet
	\begin{attlista}
		\item driva sektionens arbete med frågor som berör likabehandling, internationalisering samt arbetsmiljö
		\item erbjuda en kanal där sektionsmedlemmar kan vända sig till med frågor som berör likabehandling, internationalisering samt arbetsmiljö
		\item stötta sektionsmedlemmar som har utsatts för kränkande behandling
	\end{attlista}
	\item[Sammansättning]
	Trivselrådet består av
	\begin{itemize}
		\item Trivselmästare
		\item Likabehandlingsombud
		\item Studerandeskyddsombud
		\item Världsmästare
		\item Her Tech Future representant
	\end{itemize}

\end{reglemlista}

\section{Framtidsutskottet}

\begin{reglemlista}

	\item[Åligganden]
	Det åligger Framtidsutskottet
	\begin{attlista}
		\item bereda förslag för att presentera till styrelsen eller sektionsmötet
		\item  förankra sin arbetsplan hos styrelsen eller sektionsmötet
		\item arbeta strategiskt för att möjligöra långsiktiga mål
		\item utarbeta ett förslag till de strategiska målen till det möte definierat i stadgarna
	\end{attlista}
	\item[Sammansättning]
	Framtidsutskottet består av
	\begin{itemize}
		\item Framtidsordförande
		\item Framtidsledamöter
	\end{itemize}

\end{reglemlista}

\section{Tackmästeriet}
	\begin{reglemlista}
		\item[Åligganden]
		Det åligger Tackmästeriet
		\begin{attlista}
			\item planera och genomföra sektionsgemensamma tackevenemang
		\end{attlista}
		\item[Sammansättning]
		Tackmästeriet består av
		\begin{itemize}
			\item Tackmästare
			\item Tackmästerister
		\end{itemize}
	\end{reglemlista}

\section{Centralprocessutskottet}
	\begin{reglemlista}
		\item[Åligganden]
		Det åligger Centralprocessutskottet
		\begin{attlista}
			\item utveckla och underhålla sektionens hemsidor 
                \item utveckla och underhålla sektionens system, servrar, och vad som hör därtill
                \item tillgodose övriga programmerings- och teknikprojekt
		\end{attlista}
		\item[Sammansättning]
		Centralprocessutskottet består av
		\begin{itemize}
			\item Processmästare
                \item Vice Processmästare
                \item DWWW-ansvarig
                \item root
                \item Utvecklare
                \item Funktionärer
		\end{itemize}
	\end{reglemlista}

%%%%%%%%%%%%% SECTION %%%%%%%%%%%%%
\section{Funktionärer}

Samtliga funktionärsposter har mandatperioden kalenderår (1/1 - 31/12), om inget annat anges.

%%%%%%%%%%%%% STYRELSEPOSTER %%%%%%%%%%%%%

%%% funktionar %%%
\funktionar{Ordföranden}
Ordföranden leder styrelsens arbete och är sektionens ansikte utåt.

\begin{reglemlista}

	\item[Åligganden]
	Det åligger Ordföranden
	\begin{attlista}
		\item representera och föra sektionens talan
		\item kalla till styrelsemöten och upprätta erforderliga handlingar
		\item leda och fördela arbetet inom styrelsen
		\item kalla till sektionsmöte och upprätta erforderliga handlingar
		\item boka lokal till sektionsmöte
		\item hålla kontakt med TLTH:s ordförande samt övriga sektionsordförande
		\item hålla kontakt med Inspektor, vaktmästaren och husstyrelsen
		\item hålla kontakt med de övriga D-sektionerna i landet
		\item kontinuerligt rapportera vad som händer på sektionen via sektionens informationskanaler
		\item vid begäran skriva intyg till avgående funktionär; avgående ordförandens intyg skrivs av Inspektor tillsammans med tillträdande ordföranden
	\end{attlista}

\end{reglemlista}


%%% funktionar %%%
\funktionar{Informationsansvarig}
Informationsansvarig ansvarar för sektionens informationsspridning samt att upprätthålla
kommunikation mellan styrelsen och sektionen. Informationsansvarig leder även Informationsutskottet.

\begin{reglemlista}

	\item[Åligganden]
	Det åligger Informationsansvarig
	\begin{attlista}
		\item leda Informationsutskottet
		\item vara ansvarig utgivare av sektionens medier
		\item hålla sektionens anslagstavlor och informationsskärmar i god ordning
		\item publicera mötesprotokoll på WWW
		\item upprätthålla god kommunikation mellan styrelsen och sektionens medlemmar

		\item kontinuerligt rapportera vad som händer på sektionen via sektionens informationskanaler
	\end{attlista}

\end{reglemlista}

%%% funktionar %%%
\funktionar{Skattmästare}
Skattmästaren handhar sektionens ekonomi, sköter bokföring och ekonomisk
uppföljning. Skattmästaren leder även skattmästeriet.

\begin{reglemlista}

	\item[Åligganden]
	Det åligger Skattmästaren
	\begin{attlista}
		\item leda skattmästeriet
		\item presentera bokslut och halvårsrapport
		\item redovisa sektionens ekonomiska status för styrelsen i slutet av varje läsperiod
		\item hålla kontakten med generalsekreteraren på Teknologkåren
	\end{attlista}
	
		\item[Rättigheter]
	Skattmästaren äger rätt
	\begin{attlista}
		\item ha fortsatt tillgång till föreningens bokföring och övrig relevant dokumentation i syfte att möjliggöra upprättande och granskning av bokslut för det gångna verksamhetsåret. Denna
åtkomst ska säkerställas i minst fem månader efter det att bokslut för det föregående verksamhetsåret har fastställts av årsmötet, eller fem månader efter mandatperiodens utgång, beroende
på vilket som inträffar sist.
	\end{attlista}

\end{reglemlista}

%%% funktionar %%%
\funktionar{Studierådsordförande}
Studierådsordföranden är ansvarig för sektionens studiebevakning.
Studierådsordföranden leder även studierådet.

\begin{reglemlista}

	\item[Åligganden]
	Det åligger Studierådsordförande
	\begin{attlista}
		\item leda studierådet
		\item hålla kontakten med relevanta utbildningsnämnder
		\item aktivt deltaga i SRX, Teknologkårens studierådsordförandekollegium
	\end{attlista}

\end{reglemlista}

%%% funktionar %%%
\funktionar{Cafémästare}
Cafémästaren är ansvarig för sektionens caféverksamhet. Cafémästaren leder även cafémästeriet.

\begin{reglemlista}

	\item[Åligganden]
	Det åligger Cafémästaren
	\begin{attlista}
		\item leda cafémästeriet
	\end{attlista}

\end{reglemlista}

%%% funktionar %%%
\funktionar{Näringslivsansvarig}
Näringslivsansvarig samordnar sektionens kontakter med näringslivet.
Näringslivsansvarig leder även Näringslivsutskottet.

\begin{reglemlista}

	\item[Åligganden]
	Det åligger Näringslivsansvarig
	\begin{attlista}
		\item leda Näringslivsutskottet
		\item verka för att information om datateknik- och infocomutbildningarna vid LTH sprids till företag och näringsliv
		\item aktivt deltaga i Teknologkårens näringslivskollegium
	\end{attlista}

\end{reglemlista}

%%% funktionar %%%
\funktionar{Källarmästare}
Källarmästaren är ansvarig för sektionens lokaler och inventarier.
Källarmästaren leder även källarmästeriet.

\begin{reglemlista}

	\item[Åligganden]
	Det åligger Källarmästaren
	\begin{attlista}
		\item leda källarmästeriet
		\item administrera dörraccess och nyckelåtkomst till iDéts faciliteter

	\end{attlista}

\end{reglemlista}

%%% funktionar %%%
\funktionar{Aktivitetsansvarig}
Aktivitetsansvarig är ansvarig för sektionens studiesociala verksamheter,
undantaget fest- och pubverksamhet. Aktivitetsansvarig leder även Aktivitetsutskottet.

\begin{reglemlista}

	\item[Åligganden]
	Det åligger Aktivitetsansvarig
	\begin{attlista}
		\item leda Aktivitetsutskottet
		\item verka för att sektionen har lag med i tandemstafetten
		\item hålla kontakt med aktivitetssamordnaren på Teknologkåren samt övriga sektioners Aktivitetsansvarig
		\item anordna för sektionens medlemmar trevliga evenemang
		\item informera om större arrangemang på andra orter
	\end{attlista}

\end{reglemlista}

%%% funktionar %%%
\funktionar{Sexmästare}
Sexmästaren är ansvarig för sektionens festverksamhet. Sexmästaren leder även Sexmästeriet.

\begin{reglemlista}

	\item[Åligganden]
	Det åligger Sexmästaren
	\begin{attlista}
		\item leda sexmästeriet
	\end{attlista}

\end{reglemlista}

%%% funktionar %%%
\funktionar{Øverphøs}
Øverphøs är ansvarig för sektionens nollning.

\begin{reglemlista}

	\item[Åligganden]
	Det åligger Øverphøs
	\begin{attlista}
		\item leda Nollningsutskottet.
		\item samarbeta med Nollegeneralen, de andra Øverphøsen och LTH.
	\end{attlista}

\end{reglemlista}

%%%%%%%%%%%%% VICEPOSTER I STYRELSEN %%%%%%%%%%%%%

%%% funktionar %%%
\funktionar{Vice Ordförande}
Vice Ordförande skall vara Ordförande behjälplig, föra protokoll under styrelsemöten och sektionsmöten.

\begin{reglemlista}

	\item[Åligganden]
	Det åligger Vice Ordförande
	\begin{attlista}
		\item i Ordförandes frånvaro, leda Styrelsen.
		\item vara Ordförande behjälplig, när så denne behöver det.
		\item ansvara för att Policy-sidan, Stadgarna och Reglementet 			hålls uppdaterade.
		\item vara kontaktperson för fria föreningar och eventuella projekt
		\item anslå kallelser till styrelse- och sektionsmöten
		\item föra protokoll vid styrelsemöten och sektionsmöten samt anslå dessa
		\item arkivera uppförda protokoll och övriga handlingar i därför avsedd pärm
		\item bistå ordförande i förberedande av sektionsmöten
	\end{attlista}

	\item[Titulering]
	Alternativ titulering är Sekreterare

\end{reglemlista}

%%% funktionar %%%
\funktionar{Vice Informationsansvarig}
Vice Informationsansvarig hjälper Informationsansvarig i dennes arbete.

\begin{reglemlista}

	\item[Åligganden]
	Det åligger Vice Informationsansvarig
	\begin{attlista}
		\item hjälpa Informationsansvarig i dennes arbete
		\item överta Informationsansvarigs funktion i dennes frånvaro
	\end{attlista}

\end{reglemlista}

%%% funktionar %%%
\funktionar{Vice Skattmästare}
Vice Skattmästaren hjälper Skattmästaren i dennes arbete.

\begin{reglemlista}

	\item[Åligganden]
	Det åligger Vice Skattmästaren
	\begin{attlista}
		\item hjälpa Skattmästaren i dennes arbete.
		\item överta Skattmästarens funktion i dennes frånvaro.
		
	\end{attlista}

\end{reglemlista}

%%% funktionar %%%
\funktionar{Vice Cafémästare}
Vice Cafémästaren hjälper Cafémästaren i dennes arbete.

\begin{reglemlista}

	\item[Åligganden]
	Det åligger Vice Cafémästaren
	\begin{attlista}
		\item hjälpa Cafémästaren i dennes arbete
		\item överta Cafémästarens funktion i dennes frånvaro
	\end{attlista}

\end{reglemlista}

%%% funktionar %%%
\funktionar{Vice Näringslivsansvarig}
Vice Näringslivsansvarig hjälper Näringslivsansvarig i dennes arbete.

\begin{reglemlista}

	\item[Åligganden]
	Det åligger Vice Näringslivsansvarig
	\begin{attlista}
		\item hjälpa Näringslivsansvarig i dennes arbete
		\item i Näringslivsansvariges frånvaro leda Näringslivsutskottet
	\end{attlista}

\end{reglemlista}

%%% funktionar %%%
\funktionar{Vice Källarmästare}
Vice Källarmästaren hjälper Källarmästaren i dennes arbete.

\begin{reglemlista}

	\item[Åligganden]
	Det åligger Vice Källarmästaren
	\begin{attlista}
		\item hjälpa Källarmästaren i dennes arbete
		\item överta Källarmästarens funktion i dennes frånvaro
	\end{attlista}

\end{reglemlista}

%%% funktionar %%%
\funktionar{Vice Aktivitetsansvarig}
Vice Aktivitetsansvarig hjälper Aktivitetsansvarig i dennes arbete.

\begin{reglemlista}

	\item[Åligganden]
	Det åligger Vice Aktivitetsansvarig
	\begin{attlista}
		\item hjälpa Aktivitetsansvarig i dennes arbete
		\item överta Aktivitetsansvarigs funktion i dennes frånvaro
	\end{attlista}

\end{reglemlista}

%%% funktionar %%%
\funktionar{Vice Sexmästare}
Vice Sexmästaren hjälper Sexmästaren i dennes arbete.

\begin{reglemlista}

	\item[Åligganden]
	Det åligger Vice Sexmästaren
	\begin{attlista}
		\item hjälpa Sexmästaren i dennes arbete
		\item överta Sexmästarens funktion i dennes frånvaro
	\end{attlista}

\end{reglemlista}

%%% funktionar %%%
\funktionar{Vice Studierådsordförande}
Vice Studierådsordförande hjälper Studierådsordföranden i dennes arbete.

\begin{reglemlista}

	\item[Åligganden]
	Det åligger Vice Studierådsordföranden
	\begin{attlista}
		\item hjälpa Studierådsorföranden i dennes arbete
		\item i Studierådsordförandens frånvaro leda studierådet
	\end{attlista}
\end{reglemlista}

%%%%%%%%%%%%% ÖVRIGA FUNKTIONÄRER %%%%%%%%%%%%%

%%% funktionar %%%
\funktionar{Alumniansvarig}
\begin{reglemlista}

	\item[Åligganden]
	Det åligger alumniansvarig

	\begin{attlista}
		\item med hjälp av alumnigruppen se till att sektionens alumniverksamhet fortskrider
		\item leda alumnigruppen
	\end{attlista}

\end{reglemlista}

%%% funktionar %%%
\funktionar{Alumnigruppsmedlem}

\begin{reglemlista}

	\item[Åligganden]
	Det åligger alumnigruppsmedlem
	\begin{attlista}
		\item bistå alumniansvarig i dess arbete kring alumniverksamheten
	\end{attlista}

\end{reglemlista}



%%% funktionar %%%
\funktionar{Arkivarie}
Arkivarien ska handha sektionens arkiv.

\begin{reglemlista}

	\item[Åligganden]
	Det åligger Arkivarien
	\begin{attlista}
		\item föra och underhålla sektionens arkiv
		\item föra och underhålla sektionens dataarkiv
	\end{attlista}

\end{reglemlista}

%%% funktionar %%%
\funktionar{Artist}
Artisten utför de artistiska arbetena på sektionen.

\begin{reglemlista}

	\item[Åligganden]
	Det åligger Artisten
	\begin{attlista}
		\item leda och utföra de artistiska arbetena på sektionen
	\end{attlista}

\end{reglemlista}

%%% funktionar %%%
\funktionar{Bilansvarig}

Bilansvarig är ansvarig för sektionens bil(ar) då sådan(a) finns.

\begin{reglemlista}

	\item[Åligganden]
	Det åligger Bilansvarig
	\begin{attlista}
		\item se till att sektionens bil är i ett kördugligt skick
        \item se till att besiktning och däckbyte utförs i tid
	\end{attlista}

\end{reglemlista}


%%% funktionar %%%
\funktionar{Coach}

Coach ska för sektionen arrangera fotbollslagets verksamhet.

\begin{reglemlista}

	\item[Åligganden]
	Det åligger Coach
	\begin{attlista}
		\item hålla i fotbollslagets träningar och verksamhet.
	\end{attlista}

\end{reglemlista}


%%% funktionar %%%
\funktionar{Ljud- och ljusansvarig}
Ljud- och ljusansvarig är ansvarig för sektionens ljud- och
ljusutrustning.

\begin{reglemlista}

	\item[Åligganden]
	Det åligger Ljud- och ljusansvarig
	\begin{attlista}
		\item sköta och om serva iDéts ljud- och ljusutrustning
		\item då utrustningen skall användas åtse att kompetent personal finns på plats för att sköta densamma
	\end{attlista}

\end{reglemlista}

%%% funktionar %%%
\funktionar{DWWW-ansvarig}
DWWW-ansvarig underhåller sektionens hemsida.

\begin{reglemlista}

	\item[Åligganden]
	Det åligger DWWW-ansvarig
	\begin{attlista}
		\item agera projektledare för underhåll och utveckling av sektionens webbsidor
	\end{attlista}

\end{reglemlista}

%%% funktionar %%%
\funktionar{Redaktör}
Redaktören ska leda Redaktionen i arbetet att producera och publicera publikationer för sektionens nöje.

\begin{reglemlista}

	\item[Åligganden]
	Det åligger Redaktör
	\begin{attlista}
		\item För sektionen producera och publicera publikationer i lämpligt medium.
	\end{attlista}

\end{reglemlista}

\funktionar{Journalist}
Journalisten bistår Redaktören i dennes arbete.
\begin{reglemlista}

	\item[Åligganden]
	Det åligger Journalisten
	\begin{attlista}
		\item bistå Redaktören i dennes arbete.
	\end{attlista}

\end{reglemlista}


%%% funktionar %%%
\funktionar{Her Tech Future Representant}
Her Tech Future Representant är ansvarig för Her Tech Future-arrangemanget på
sektionen.

\begin{reglemlista}

	\item[Åligganden]
	Det åligger Her Tech Future Representant
	\begin{attlista}
		\item framföra en positiv bild av LTH
		\item ansvara för att sektionen är representerad på ett bra sätt under Her Tech Future
	\end{attlista}

	\item[Mandatperiod]
	Her Tech Future Representant mandatperiod är 1 januari - 31 december

\end{reglemlista}

%%% funktionar %%%
\funktionar{Fotograf}
Fotografen ska dokumentera händelser på sektionen för eftervärlden.

\begin{reglemlista}

	\item[Åligganden]
	Det åligger Fotografen
	\begin{attlista}
		\item fotografera alla för sektionen och mänskligheten viktiga händelser på sektionen
		\item snarast möjligt göra fotografierna tillgängliga för sektionens medlemmar
	\end{attlista}

\end{reglemlista}

%%% funktionar %%%
\funktionar{Filmare}
Filmaren dokumenterar och redigerar i videoformat händelser på sektionen.

\begin{reglemlista}

	\item[Åligganden]
	Det åligger Filmaren
	\begin{attlista}
		\item filma alla för sektionen och mänskligheten viktiga händelser på sektionen
		\item vid behov redigera inspelat material för publicering
		\item vid lämplig tidpunkt göra filmerna tillgängliga för sektionens medlemmar
	\end{attlista}

\end{reglemlista}

%%% funktionar %%%
\funktionar{Hovmästare}
Hovmästaren är den person som har det övergripande ansvaret för sexmästeriets servering och dukning.

\begin{reglemlista}

	\item[Åligganden]
	Det åligger Hovmästaren
	\begin{attlista}
		\item vara ansvarig för Sexmästeriets servering, dukning, bordsplacering och schemaläggning.
		\item tillgodose sektionens behov av sittningsdryck till sittningar.
	\end{attlista}

\end{reglemlista}

%%% funktionar %%%
\funktionar{Vinförman}
Vinförmannen är den person som har det övergripande ansvaret för sexmästeriets vinsortiment.

\begin{reglemlista}

	\item[Åligganden]
	Det åligger Vinförmannen
	\begin{attlista}
		\item i samråd med pubmästaren och barmästaren se till att sexmästeriet har ett tillräckligt vinsortiment till pubar och andra tillställningar exklusive sittningar
		\item hjälpa hovmästaren att tillgodose sektionen med sittningsdryck till sittningar

	\end{attlista}

\end{reglemlista}

%%% funktionar %%%
\funktionar{Idrottsförman}
Idrottsförmannen ska tillse sektionsmedlemmarnas fysiska hälsa.

\begin{reglemlista}

	\item[Åligganden]
	Det åligger Idrottsförmannen
	\begin{attlista}
		\item tillgodose medlemmarnas behov av motion
		\item gå på kollegiemöte för kårens idrottsförmän när så 		sammankallas
	\end{attlista}

\end{reglemlista}

%%% funktionar %%%
\funktionar{D-sportare}
D-sportaren arrangerar för sektionen turneringar inom E-sport

\begin{reglemlista}

	\item[Åligganden]
	Det åligger Idrottsförmannen
	\begin{attlista}
		\item att arrangera E-sportsturneringar för sektionen
	\end{attlista}

\end{reglemlista}

%%% funktionar %%%
\funktionar{Jubileumsansvarig}
Jubileumsansvarig arrangerar, under Jubileumsgeneralens ledning, D-sektionensjubileum.

\begin{reglemlista}
	\item[Åligganden]
	Det åligger Jubileumsansvarig
	\begin{attlista}
		\item bistå Jubileumsgeneralens arbete under jubileumet
	\end{attlista}
	
	\item[Mandatperiod]
	Jubileumsansvarigs mandatperiod är 1 juli - 30 juni
\end{reglemlista}

%%% funktionar %%%
\funktionar{Jubileumsgeneral}
Jubileumsgeneral ansvarar för arrangerandet av D-sektionens jubileum.

\begin{reglemlista}
	\item[Åligganden]
	Det åligger Jubileumsgeneralen
	\begin{attlista}
		\item ha det övergripande ansvaret för arrangerandet av D-sektionens jubileum
		\item informera sektionens styrelse om jubileumets fortskridande
		\item före jubileumets start presentera budgetförslag för sektionens styrelse
	\end{attlista}
	\item[Mandatperiod]
	Jubileumsgeneralens mandatperiod är 1 juli - 30 juni
\end{reglemlista}

%%% funktionar %%%
\funktionar{Trivselmästare}
Trivselmästaren är ansvarig för likabehandling- och internationaliseringsarbetet på sektionen. Trivselmästaren ska driva sektionens arbetsmiljöfrågor. Trivselmästaren leder trivselrådet. Trivselmästaren bör
ha tystnadsplikt för ärenden den behandlar.

\begin{reglemlista}

	\item[Åligganden]
	Det åligger Trivselmästaren

	\begin{attlista}
		\item tillsammans med Likabehandlingsombuden driva sektionens likabehandlingsarbete
		\item ansvara för sektionens internationalisering
		\item tillsammans med Studerandeskyddsombud, bevaka Sektionens intressen i arbetsmiljöfrågor
		\item informera styrelsen om arbetet som sker inom trivselrådet
	\end{attlista}

\end{reglemlista}

%%% funktionar %%%
\funktionar{Likabehandlingsombud}
Likabehandlingsombudet bör ha tystnadsplikt för ärenden den behandlar.

\begin{reglemlista}

	\item[Åligganden]
	Det åligger Likabehandlingsombudet

	\begin{attlista}
		\item bevaka sektionens verksamhet från ett likabehandlingsperspektiv
		\item uppmärksamma styrelsen på situationer och miljöer som skulle kunna upplevas som kränkande av medlemmar ur sektionen eller bryta mot Teknologkårens likabehandlingspolicy
		\item driva likabehandlingsfrågor inom utbildningen i Trivselrådet
	\end{attlista}


\end{reglemlista}

%%% funktionar %%%
\funktionar{Karnevalsansvarig}
Karnevalsansvarig har hand om sektionens deltagande i karnevalen och väljs var fjärde år.

\begin{reglemlista}

	\item[Åligganden]
	Det åligger Karnevalsansvarig
	\begin{attlista}
		\item se till att sektionen deltager i karnevalens arrangemang såsom tåg- och tältnöjen
		\item presentera en budget över sin verksamhet för styrelsen senast
		tre månader före karnevalen äger rum
		\item hålla kontakten med Karnevalskommittén
	\end{attlista}

	\item[Mandatperiod]
	Karnevalsansvariges mandatperiod är 1 juli - 30 juni

	\item[Val]
	Karnevalsansvarig väljs var fjärde år

\end{reglemlista}

%%% funktionar %%%
\funktionar{Köksmästare}
Köksmästaren har det övergripande ansvaret för Sexmästeriets köksverksamhet under deras evenemang.

\begin{reglemlista}

	\item[Åligganden]
	Det åligger Köksmästaren
	\begin{attlista}
		\item vara ansvarig för Sexmästeriets matlagning under deras evenemang
	\end{attlista}
\end{reglemlista}

%%% funktionar %%%
\funktionar{Sektionskock}
\begin{reglemlista}
	\item[Åligganden]
	Det åligger Sektionskocken
	\begin{attlista}
		\item laga mat utanför sittningar och pubar.
	\end{attlista}
\end{reglemlista}

%%% funktionar %%%
\funktionar{Preferensmästare}
\begin{reglemlista}
	\item[Åligganden]
	Det åligger Preferensmästaren
	\begin{attlista}
		\item vara ansvarig för Sexmästeriets köksverksamhet gällande specialkost under deras evenemang.
	\end{attlista}
\end{reglemlista}

%%% funktionar %%
\funktionar{Vice Köksmästare}
Vice Köksmästaren hjälper Köksmästaren i dennes arbete.
\begin{reglemlista}

	\item[Åligganden]
	Det åligger Vice Köksmästaren
	\begin{attlista}
		\item hjälpa Köksmästaren i dennes arbete
		\item överta Köksmästarens funktion vid dennes frånvaro
	\end{attlista}

\end{reglemlista}

%%% funktionar %%%
\funktionar{LAN-partyansvarig}
LAN-partyansvarig ansvarar för att anordna LAN-partyn.

\begin{reglemlista}

	\item[Åligganden]
	Det åligger LAN-partyansvarig
	\begin{attlista}
		\item anordna LAN-partyn
	\end{attlista}

\end{reglemlista}


%%% funktionar %%%
\funktionar{Medaljelelekommittémedlem}
Medaljelelekommittémedlemen ingår i medaljelelekommittén och sköter om utdelning av medaljer och riddarskap.

\begin{reglemlista}

	\item[Åligganden]
	Det åligger Medaljelelekommittémedlemmen
	\begin{attlista}
		\item närvara på Medaljelelekommitténs möten
		\item lusteligen dela ut medaljer och andra utmärkelser till funktionärer som gjort sig förtjänta därav
	\end{attlista}
\end{reglemlista}

%%% funktionar %%%
\funktionar{Medlem i Projektgruppen för Teknikfokus}

\begin{reglemlista}
	\item[Åligganden]
	Det åligger Medlem i Projektgruppen för Teknikfokus
	\begin{attlista}
		\item bistå Teknikfokusansvarig i dennes arbete
	\end{attlista}

	\item[Mandatperiod]
	Mandatperioden för Medlem i Projektgruppen för Teknikfokus är 1 juli - 30 juni.
\end{reglemlista}

%%% funktionar %%%
\funktionar{Mentorsansvarig}
Mentorsanvarig är ytterst ansvarig för all verksamhet som rör Mentorsprogrammet.

\begin{reglemlista}

	\item[Åligganden]
	Det åligger Mentorsansvarig
	\begin{attlista}
		\item driva och utveckla D-sektionens Mentorsprogram
	\end{attlista}

\end{reglemlista}

%%% funktionar %%%
\funktionar{Märkvärdig}

Märkvärdig ska hjälpa Shopaholic med D-shopens inventarie. Detta bland annat genom att ge design för ny merch och nya märken.

\begin{reglemlista}

	\item[Åligganden]
	Det åligger Märkvärdig
	\begin{attlista}
		\item bistå Shopaholic med D-shopens inventarie
		\item  ge förslag på design på merch eller märken som kan säljas i D-shopen
	\end{attlista}

\end{reglemlista}


%%% funktionar %%%
\funktionar{Nöjesförman}

Nöjesförmannen ska för sektionen arrangera evenemang av lek-, mys- och pysselkaraktär.

\begin{reglemlista}

	\item[Åligganden]
	Det åligger Nöjesförman
	\begin{attlista}
		\item för sektionen arrangera evenemang av lek-, mys- och pysselkaraktär
	\end{attlista}

\end{reglemlista}


%%% funktionar %%%
\funktionar{Shopaholic}

Shopaholic ska leda D-shopen, ansvara för dess inventarie och tillse sektionen med märken, merch, sångböcker och ouvvar.

\begin{reglemlista}

	\item[Åligganden]
	Det åligger Shopaholic
	\begin{attlista}
		\item ansvara för försäljning, inköp och inventering av D-shopens produkter
		\item sköta kontakten med andra utskott där D-shopens inventarie säljs
	\end{attlista}

\end{reglemlista}

%%% funktionar %%%
\funktionar{Skald}

Skald ansvarar för sångboken och ser till att den är uppdaterad och studentikos.

\begin{reglemlista}

	\item[Åligganden]
	Det åligger Skald
	\begin{attlista}
		\item se över och revidera sångboken
		\item att bistå Shopaholic med D-shopens inventarie samt försäljning
	\end{attlista}

\end{reglemlista}

%%% funktionar %%
\funktionar{Phadder}
Phaddrar verkar under nollningen för att välkomna de nya studenterna till D-sektionen.

\begin{reglemlista}

	\item[Åligganden]
	Det åligger Phadder
	\begin{attlista}
		\item vara Nollningsutskottet behjälplig under nollningen
		\item välkommna de nya studenterna till D-sektionen
	\end{attlista}

	\item[Mandatperiod]
	Phadderns mandatperiod är 1~april -- 31~oktober

\end{reglemlista}

\funktionar{Pluggphadder}
Pluggphaddrar verkar under nollningen för att hjälpa de nya studenterna komma igång med studierna.

\begin{reglemlista}

	\item[Åligganden]
	Det åligger Pluggphadder
	\begin{attlista}
		\item Pluggphaddrar ska agera studiehjälp till nya studenter under Nollningen
		\item på ett sunt vis introducera de nya studenterna till dess studier på D-sektionen
        \item vara Studierådet behjälplig under nollningen
	\end{attlista}

	\item[Mandatperiod]
	Pluggphadderns mandatperiod är 1~april -- 31~oktober

\end{reglemlista}

%%% funktionar %%%
\funktionar{Pubmästare}
Pubmästaren ansvarar för pubverksamheten.

\begin{reglemlista}

	\item[Åligganden]
	Det åligger Pubmästaren
	\begin{attlista}
		\item ansvara för sexmästeriets pubverksamhet.
		\item tillsammans med barmästaren hålla ett högkvalitativt öl- och drinksortiment samt se till att flera alkoholfria alternativ alltid finns.
	\end{attlista}

\end{reglemlista}

%%% funktionar %%%
\funktionar{Vice Pubmästare}
Vice Pubmästaren hjälper Pubmästaren i dennes arbete.

\begin{reglemlista}

	\item[Åligganden]
	Det åligger Vice Pubmästaren
	\begin{attlista}
		\item hjälpa Pubmästaren i dennes arbete
		\item överta Pubmästarens funktion vid dennes frånvaro
	\end{attlista}

\end{reglemlista}

%%% funktionar %%%
\funktionar{Barmästare}
Barmästaren ansvarar för barverksamheten under sittningar.

\begin{reglemlista}

	\item[Åligganden]
	Det åligger Barmästaren
	\begin{attlista}
		\item anordna barverksamhet i samband med sittningar och eftersläpp arrangerade av sexmästeriet
		\item tillsammans med pubmästaren hålla ett högkvalitativt öl- och drinksortiment samt se till att flera alkoholfria alternativ alltid finns
	\end{attlista}

\end{reglemlista}

%%% funktionar %%%
\funktionar{Vice Barmästare}
Vice Barmästaren hjälper Barmästaren i dennes arbete.

\begin{reglemlista}

	\item[Åligganden]
	Det åligger Vice Barmästaren
	\begin{attlista}
		\item hjälpa Barmästaren i dennes arbete
		\item överta Barmästarens roll vid dennes frånvaro
	\end{attlista}

\end{reglemlista}


%%% funktionar %%%
\funktionar{Revisor}
Revisorerna skall granska sektionens verksamhet och bokslut.

\begin{reglemlista}

	\item[Krav]
	Det krävs
	\begin{attlista}
		\item revisorerna är myndiga svenska medborgare
		\item revisorerna har den insikt i ekonomiska förhållanden som uppdraget kräver
		\item en av revisorerna har god insyn i sektionens verksamhet
		\item revisorerna inte är försatta i konkurs eller har näringsförbud
		\item revisorerna inte har förvaltare
		\item revisorerna inte är styrelseledamöter
		\item revisorerna bör inte vara gift, sambo eller nära släkt med någon i styrelsen
	\end{attlista}

	\item[Åligganden]
	Det åligger revisorerna
	\begin{attlista}
		\item granska sektionens böcker och räkenskaper
		\item taga del av sektionsmötenas och styrelsens protokoll
		\item verkställa fortlöpande inventering eller kontrollera ställd inventering av sektionens kassa och övriga tillgångar
		\item tillse huruvida sektionens organisation av och kontroll över bokföring är tillfredsställande
		\item senast fjorton dagar innan vårterminsmötet, året
		efter mandatperioden, till styrelsen inlämna revisionsberättelse
	\end{attlista}

	\item[Rättigheter]
	Revisorerna äger rätt
	\begin{attlista}
		\item närhelst han eller hon önskar taga del av sektionens samtliga räkenskaper, protokoll och andra handlingar
		\item beredas tillträde till alla sektionens lokaler
		\item erhålla upplysningar rörande sektionens verksamhet och förvaltning
		\item närvara vid sektionens samtliga möten som ständigt adjungerade
	\end{attlista}

	\item[Skyldigheter]
	Revisorerna är skyldiga
	\begin{attlista}
		\item alltid handla för sektionens bästa
	\end{attlista}

\end{reglemlista}

%%% funktionar %%%
\funktionar{root}
root underhåller sektionens datorer och därtill kopplad utrustning.
\begin{reglemlista}

	\item[Åligganden]
	Det åligger root
	\begin{attlista}
		\item agera projektledare för underhåll och förbättrande av sektionens datorer och därtill kopplad utrustning
	\end{attlista}

\end{reglemlista}

%%% funktionar %%%
\funktionar{Studerandeskyddsombud}
Studerandeskyddsombudet bevakar sektionens intressen i olika organ för arbetsmiljöfrågor.
Studerandeskyddsombudet för fram sektionens och studierådets åsikter gentemot E-husets
husstyrelse. Studerandeskyddsombudet ansvarar för att informera studierådet om det som tagits upp på mötet för husstyrelsen.

\begin{reglemlista}

	\item[Åligganden]
	Det åligger Studerandeskyddsombudet
	\begin{attlista}
		\item verka för god studiemiljö
		\item medverka på HMS-kommitténs möten
		\item att medverka på husstyrelsemöten för E-husets husstyrelse som
		studentrepresentant
		\item förmedla sektionens åsikter på HMS-kommitémöten och husstyrelsemöten
		\item informera sektionen om kommande HMS-kommitémöten och Husstyrelsemöten
		\item informera studierådet om vad som sagts på tidigare HMS-kommitémöten och Husstyrelsemöten
		\item verka som Brandskyddssamordnare för D-sektionen
		
	\end{attlista}
	
	\item[Mandatperiod] 
	Studentrepresentanter inom Husstyrelse har mandatperioden 1~juli -- 30~juni

\end{reglemlista}

%%funktionar %%
\funktionar{Stekare}
Stekaren tillagar varor som används i maträtterna som serveras i Café iDét.
\begin{reglemlista}
	\item[Åligganden]
	Det åligger Stekaren
	\begin{attlista}
		\item med lämpligt intervall
		tillaga mat åt Cafémästeriets dagliga verksamhet, så att det alltid finns tillgängligt
	\end{attlista}

\end{reglemlista}

%%funktionar %%
\funktionar{Dagsansvarig}

\begin{reglemlista}
	\item[Åligganden]
	Det åligger Dagsansvarig
	\begin{attlista}
		\item ansvara för sektionens café under dennes tilldelade arbetspass.
	\end{attlista}

\end{reglemlista}

%%funktionar %%
\funktionar{Brunchmästare}

\begin{reglemlista}
	\item[Åligganden]
	Det åligger Brunchmästare
	\begin{attlista}
		\item planera och genomföra bruncher eller andra evenemang för Cafémästeriets räkning
        \item planera temaenligt sortiment inför temadagar
        \item koordinera beredning av bakelser som kan säljas vid passande tillfällen
	\end{attlista}

\end{reglemlista}

\funktionar{Bakis}
Bakis hjälper brunchmästaren i dennes arbete.

\begin{reglemlista}
	\item[Åligganden]
	Det åligger Bakis
	\begin{attlista}
		\item hjälpa brunchmästaren i dennes arbete
	\end{attlista}

\end{reglemlista}

\funktionar{Inventarieansvarig}

\begin{reglemlista}
	\item[Åligganden]
	Det åligger Inventarieansvarig
	\begin{attlista}
		\item ha övergripande ansvar över sortimentet i cafét
        \item se till så att inköp av ingredienser och sortiment till cafét sker
        \item i samråd med mästarna, utskottet och sektionen utveckla caféts sortiment
        \item ansvara för sektionens läskinventarier
	\end{attlista}

\end{reglemlista}

%%% funktionar %%%
\funktionar{Stabsmedlem}
Stabsmedlem arrangerar, under Øverphøs ledning, nollning.

\begin{reglemlista}

	\item[Åligganden]
	Det åligger Stabsmedlem
	\begin{attlista}
		\item bistå Øverphøs i dennes arbete.
		\item att planera och genomföra nollningsrelaterade
evenemang.
	\end{attlista}

\end{reglemlista}


%%% funktionar %%%
\funktionar{Studentrepresentant, Institutionsstyrelse}
Studentrepresentanten för fram sektionens och studierådets åsikter gentemot institutionsstyrelsen som studenten sitter i. Studentrepresentanten ska förmedla sektionens åsikter inför, under och efter dessa möten. Åsikterna som förmedlas ska vara förankrade i studierådet. Studentrepresentanten ansvarar för att informera studierådet om det som tagits upp på institutionsstyrelsemötena.

\begin{reglemlista}

	\item[Åligganden]
	Det åligger Studentrepresentant, Institutionsstyrelse
	\begin{attlista}
		\item medverka på institutionsstyrelsemötena för den institution som studentrepresentanten blivit vald till av TLTHs fullmäktige
		\item förmedla studierådets åsikter på institutionsstyrelsemötena
		\item informera studierådet om kommande Institutionsstyrelsemöten
		\item informera studierådet om vad som sagts på tidigare Institutionsstyrelsemöten
	\end{attlista}
        \item[Mandatperiod]
        Studentrepresentanter inom Institutionsstyrelse har mandatperioden 1 juli - 30 juni

\end{reglemlista}

%%% funktionar %%%
\funktionar{Studentrepresentant, Programledning}
Studentrepresentanten för fram sektionens och studierådets åsikter gentemot den programledning som studenten sitter i. Studentrepresentanten ska förmedla sektionens åsikter inför, under och efter dessa möten. Åsikterna som förmedlas ska vara förankrade i studierådet. Studentrepresentanten ansvarar för att informera studierådet om det som tagits upp på mötet för programledningen

\begin{reglemlista}

	\item[Åligganden]
	Det åligger Studentrepresentant, Programledning
	\begin{attlista}
		\item medverka på programledningsmötena för den programledning som studentrepresentanten blivit vald till av TLTHs fullmäktige
		\item förmedla studierådets åsikter på programledningsmötena
		\item informera studierådet om kommande Programledningsmöten
		\item informera studierådet om vad som sagts på tidigare Programledningsmöten
	\end{attlista}
         \item[Mandatperiod]
        Studentrepresentanter inom Programledningen har mandatperioden 1 juli - 30 juni
\end{reglemlista}


%%% funktionar %%%
\funktionar{Studierådssekreterare}
Studierådssekreterare för protokoll under studierådetsmöten.

\begin{reglemlista}

	\item[Åligganden]
	Det åligger Studierådssekreteraren
	\begin{attlista}
		\item föra protokoll under studierådsmöten
		\item se till att dessa förs upp på D-sektionens hemsida
       \end{attlista}
       
        \item[Mandatperiod]
         Studierådssekreterare har mandatperioden 1 juli - 30 juni

	\item[Val]
	Studierådssekreteraren väljs av Studierådet

\end{reglemlista}

%%% funktionar %%%
\funktionar{Sångarstridsförman}
Sångarstridsförmannen är ansvarig för sektionens deltagande i Sångarstriden.

\begin{reglemlista}

	\item[Åligganden]
	Det åligger Sångarstridsförmannen
	\begin{attlista}
		\item se till att sektionen ställer upp i Sångarstriden
	\end{attlista}

\end{reglemlista}

%%% funktionar %%%
\funktionar{Sångförman}
Sångförmannen leder sången under sektionens fester.

\begin{reglemlista}

	\item[Åligganden]
	Det åligger Sångförmannen
	\begin{attlista}
		\item ansvara för schema och sång under sektionens sittningar.
		\item  hjälpa Hovmästaren sätta tidschema på
sektionens sittningar.
		\item att vara sexmästeriets kontaktperson angående gyckel.
		\item lära nollan vett och etikett på sittningar.
		\item göra sångblad då så erfodras.
		\item lära nollan nollevisan.
	\end{attlista}

\end{reglemlista}

%%% funktionar %%%
\funktionar{Talman}
Talmannen är mötesordförande på sektionsmötena.

\begin{reglemlista}

	\item[Åligganden]
	Det åligger Talmannen
	\begin{attlista}
		\item sitta som mötesordförande under sektionsmöten
		\item hjälpa styrelsen med förberedelserna inför sektionsmöten
	\end{attlista}

\end{reglemlista}

%%% funktionar %%%
\funktionar{Tandemgeneral}
Tandemgeneralen ska representera D-sektionen på Tandemstafetten, delta
vid dess planering och se till att D-sektionens medlemmar får chansen att delta.

\begin{reglemlista}

	\item[Åligganden]
	Det åligger Tandemgeneralen
	\begin{attlista}
		\item ge D-sektionens medlemmar en chans att delta i Tandemstafetten
	\end{attlista}

\end{reglemlista}

%%% funktionar %%%
\funktionar{Teknikfokusansvarig}
Teknikfokusansvarig har tillsammans med Näringslivsansvarig hand om
D-sektionens årliga arbetsmarknadsmässa Teknikfokus.

\begin{reglemlista}

	\item[Åligganden]
	Det åligger Teknikfokusansvarig
	\begin{attlista}
		\item tillsammans med Näringslivsansvarig ansvara för arrangerandet av Teknikfokus
		\item hjälpa Näringslivsansvarig att i god tid presentera en budget för Teknikfokus
	\end{attlista}
	\item[Mandatperiod]
	Teknikfokusansvariges mandatperiod är 1 juli - 30 juni
\end{reglemlista}

%%% funktionar %%%
\funktionar{Trädgårdsmästare}
Trädgårdsmästaren har hand om sektionens växter.

\begin{reglemlista}

	\item[Åligganden]
	Det åligger Trädgårdsmästaren
	\begin{attlista}
		\item sköta sektionens växter
	\end{attlista}

\end{reglemlista}

%%% funktionar %%%
\funktionar{TappaD}
TappaD ansvarar för sektionens tappsystem

\begin{reglemlista}

	\item[Åligganden]
	Det åligger TappaD
	\begin{attlista}
		\item ansvara för underhåll och kunskapsbevaring av sektionens tappsystem
	\end{attlista}

\end{reglemlista}

%%% funktionar %%%
\funktionar{Utedischoansvarig}
Utedischoansvarig ansvarar för D-sektionens arrangemang i det årliga utedischot.

\begin{reglemlista}

	\item[Åligganden]
	Det åligger Utedischoansvarig
	\begin{attlista}
		\item arrangera utedischo i samband med nollningen
	\end{attlista}
\end{reglemlista}

%%% funktionar %%%
\funktionar{Valberedningens ordförande}
Valberedningens ordförande leder valberedningens arbete.

\begin{reglemlista}

	\item[Åligganden]
	Det åligger Valberedningens ordförande
	\begin{attlista}
		\item leda valberedningens arbete
		\item inte diskutera intervjuerna och detaljer om kandidaterna förutom med valberedningen och den berörda kandidaten
\end{attlista}

	\item[Mandatperiod]
	Valberedningens ordförandens mandatperiod är 1~juli -- 30~juni.

\end{reglemlista}

%%% funktionar %%%
\funktionar{Valberedningsrepresentant}
Valberedningsrepresentanten verkar för att bredda sektionens
rekryteringsbas.

\begin{reglemlista}

	\item[Åligganden]
	Det åligger Valberedningsrepresentanten
	\begin{attlista}
		\item bredda sektionens rekryteringsbas
		\item inte diskutera intervjuerna och detaljer om kandidaterna förutom med valberedningen och den berörda kandidaten
\end{attlista}

	\item[Mandatperiod]
	Valberedningsrepresentantens mandatperiod är 1~juli -- 30~juni.

\end{reglemlista}

%%% funktionar %%%
\funktionar{Valnämndsrepresentant TLTH}
Valnämndsrepresentanten TLTH är sektionens representant i TLTHs valnämnd.

\begin{reglemlista}

	\item[Åligganden]
	Det åligger Valnämndsrepresentanten TLTH
	\begin{attlista}
		\item representera sektionen i TLTHs valnämnd
	\end{attlista}

\end{reglemlista}

%%% funktionar %%%
\funktionar{Världsmästare}
Världsmästaren verkar för sektionens internationalisering.

\begin{reglemlista}

	\item[Åligganden]
	Det åligger Världsmästaren
	\begin{attlista}
		\item representera sektionen i TLTHs internationella utskott
		\item verka för sektionens internationalisering
		\item ge internationella studenter den information de behöver
		\item öka medvetenhet om internationella studenter på sektionen
	\end{attlista}

\end{reglemlista}

%%% funktionar %%%
\funktionar{Årskursrepresentant}
Årskursrepresentanten är den som utvärderar våra kurser och framför sina
kursarers åsikter.

\begin{reglemlista}

	\item[Åligganden]
	Det åligger Årskursrepresentanten
	\begin{attlista}
		\item agera representant för sina kursare inför kursansvarig och framföra deras åsikter
		\item delta på årskursens CEQ-möten
		\item skriva slutkommentarer för kursen
        \item informera Studierådet om vad som sagts på halvtidsmöten och CEQ-möten”
	\end{attlista}
	\item[Mandatperiod]
	Årskursrepresentants mandatperiod är 1~juli -- 30~juni.

	\item[Val]
	Årskursrepresentant väljs av Studierådet.
\end{reglemlista}


%%% funktionar %%%
\funktionar{Ölförman}
Ölförmannen är i samråd med pub- och barmästare ansvarig för sexmästeriets inköp av bryggerivaror.

\begin{reglemlista}

	\item[Åligganden]
	Det åligger Ölförmannen
	\begin{attlista}
		\item i samråd med pubmästaren och barmästaren se till så att sexmästeriet har nödvändiga drycker exklusive sittningsdryck i god tid före varje tillställning
	\end{attlista}

\end{reglemlista}

%%% funktionar %%%
\funktionar{Övermarskalk}
Övermarskalken bär sektionens fana samt sköter inbjudningar vid större sektionsfester.

\begin{reglemlista}

	\item[Åligganden]
	Det åligger Övermarskalken
	\begin{attlista}
		\item vid högtider såsom Regatta och Kårbal bära sektionens fana
		\item handha medaljer och av sektionen fastställda utmärkelser samt utdelandet av dessa
		\item övervaka Ordenssällskapet D-sektionen inom TLTH och tillse att det verkar enligt gällande reglemente och stadgar
		\item handha inbjudningar och anmodningar vid sektionens festligheter
		\item ansvara för inhandling och förvaltning av ordensband
		\item ha hand om sektionsflaggan
	\end{attlista}

\end{reglemlista}

%%% funktionar %%%
\funktionar{Semesterfirare}

\begin{reglemlista}

	\item[Åligganden]
	Det åligger Semesterfiraren
	\begin{attlista}
		\item se till att sektionens medlemmar får åka på utomlandssemester till andra universitet i Sverige
	\end{attlista}

\end{reglemlista}

%%% funktionar %%%
\funktionar{\O verpeppare}
\O verpepparna agerar Stabens högra hand genom att vara med i planering
och genomförande av nollningsrelaterade evenemang. \O verpepparna ansvarar för att leda Pepparna.
\begin{reglemlista}
	\item[Åligganden]
	Det åligger \O verpeppare
	\begin{attlista}
		\item planera och genomföra nollningsrelaterade evenemang.
		\item leda Pepparna
	\end{attlista}

\end{reglemlista}

%%% funktionar %%%
\funktionar{Peppare}
Pepparna agerar \O verpepparnas högra hand genom att vara med i planering
och genomförande av nollningsrelaterade evenemang.
\begin{reglemlista}
	\item[Åligganden]
	Det åligger Pepparna
	\begin{attlista}
		\item planera och genomföra nollningsrelaterade evenemang.
	\end{attlista}

\end{reglemlista}

%%% funktionar %%%
\funktionar{Framtidsordförande}
Framtidsordförande ansvarar för att leda Framtidsutskottet.
\begin{reglemlista}
	\item[Åligganden]
	Det åligger Framtidsordförande
	\begin{attlista}
		\item kalla till möten för Framtidsutskottet och upprätta erforderliga handlingar
		\item leda och fördela arbetet inom utskottet
		\item samverka med styrelsen för att gynna sektionens strategiska arbete
	\end{attlista}

\end{reglemlista}
  
%%% funktionar %%%
\funktionar{Framtidsledamot}
Framtidsledamoten sitter i Framtidsutskottet
\begin{reglemlista}
	\item[Åligganden]
	Det åligger Framtidsledamöter
	\begin{attlista}
		\item arbeta med sektionens strategiska arbete
		\item bereda förslag för att presentera till styrelsen eller sektionsmötet
		\item bistå Framtidsordförannde i dess arbete
	\end{attlista}

\end{reglemlista}
  
%%% funktionar %%%
\funktionar{Skattförman}
Skattförmannen hjälper Skattmästaren i dennes arbete.
\begin{reglemlista}
	\item[Åligganden]
	Det åligger Skattförmannen
	\begin{attlista}
		\item att hjälpa Skattmästaren i dennes arbete.
	\end{attlista}

\end{reglemlista}

%%% funktionar %%%
\funktionar{Tackmästare}
Tackmästaren ansvarar för att leda Tackmästeriet.
\begin{reglemlista}
    \item[Åligganden]
      Det åligger Tackmästaren
      \begin{attlista}
        \item välja in Tackmästerister
        \item leda och fördela arbetet inom utskottet
        \item tillsammans med utskottet planera och genomföra sektionsgemensamma tackevenemang
        \item presentera en plan och en budget för tackevenemanget till styrelsen
      \end{attlista}
    \item[Rättigheter]
      Tackmästaren äger rätt
      \begin{attlista}
        \item delta på nästkommande sektionsgemensamma tackevenemang efter avslutad mandatperiod
      \end{attlista}
    \item[Mandatperiod]
      Tackmästarens mandatperiod är
      \begin{itemize}
      \item 1~januari -- 30~juni
      \item 1~juli -- 31~december
      \end{itemize}
\end{reglemlista}

%%% funktionar %%%
\funktionar{Tackmästerist}
\begin{reglemlista}
\item[Åligganden]
  Det åligger Tackmästeristen
  \begin{attlista}
    \item planera och genomföra sektionstacken
    \end{attlista}
  \item[Rättigheter]
    Tackmästeristen äger rätt
    \begin{attlista}
      \item delta på nästkommande sektionsgemensamma tackevenemang efter avslutad mandatperiod
      \end{attlista}
    \item[Mandatperiod]
      Tackmästeristens mandatperiod är
      \begin{itemize}
      \item 1~januari -- 30~juni
      \item 1~juli -- 31~december
      \end{itemize}
\end{reglemlista}

%%% funktionar %%%
\funktionar{Processmästare}
Processmästaren ansvarar för att leda Centralprocessutskottet och administrera sektionens hemsida.
\begin{reglemlista}
	\item[Åligganden]
        Det åligger Processmästaren	
        \begin{attlista}
		\item leda och fördela arbetet inom utskottet
            \item tillhandahålla administrativ information på hemsidan, såsom att administrera poster på sektionen
            \item hjälpa våra medlemmar att använda hemsidan
	\end{attlista}
\end{reglemlista}

%%% funktionar %%%
\funktionar{Vice Processmästare}
Vice Processmästaren hjälper Processmästaren i dennes arbete.
\begin{reglemlista}
	\item[Åligganden]
        Det åligger Vice Processmästaren	
        \begin{attlista}
		\item hjälpa DWWW-ansvarig och root i deras arbete
            \item att vägleda funktionärer och övriga sektionsmedlemmar i projekt och uppgifter inom Centralprocessutskottet 
	\end{attlista}
\end{reglemlista}

%%% funktionar %%%
\funktionar{Utvecklare}
Utvecklarna hjälper utskottet i dess åligganden.
\begin{reglemlista}
	\item[Åligganden]
        Det åligger Utvecklarna	
        \begin{attlista}
		\item hjälpa DWWW-ansvarig och root i deras arbete
            \item vägleda funktionärer och övriga sektionsmedlemmar i projekt och uppgifter inom  Centralprocessutskottet
	\end{attlista}
\end{reglemlista}

%%%%%%%%%%%%% SECTION %%%%%%%%%%%%%
\section{Medaljer och Utmärkelser}

I detta avsnitt listas de medaljer och utmärkelser som sektionen delar ut samt de kriterier som ligger bakom förvärvandet av dessa.


Medaljer och utmärkelser delas företrädesvis ut två gånger om året, under de årliga ordenskapitlen i samband med Skiphtesgasque och Nollegasque. Vid speciella tillfällen, såsom t.ex. Jubileum, kan ett extra ordenskapitel
hållas.

\subsection{Funktionärsmedaljen}
Funktionärsmedaljen erhålles av funktionärer efter avslutad mandatperiod. Denna medalj bäres med stolthet. Funktionärsmedaljen utdelas till en och samma person endast en gång.

\subsection{Stor D-medalj}
Stor D-medalj utdelas av Medaljelelekommittén på årligt ordenskapitel, till person som Medaljelelekommittén har funnit dels under en ansenlig mängd år innehaft avancerade förtroendeuppdrag inom D-sektionen, dels verkat på det mest föredömliga sätt samt åstadkommit exceptionella fördelar för D-sektionen. Stor D-medalj utdelas till en och samma person endast en gång.

\subsection{D-medalj}
D-medalj utdelas av Medaljelelekommittén på årligt ordenskapitel, till person som Medaljelelekommittén har funnit dels under flera år innehaft förtroendeuppdrag inom D-sektionen, dels på olika sätt verkat mycket föredömligt för D-sektionen. I folkmun kallat för ''riddare''. D-medalj utdelas till en och samma person endast en gång.

\subsection{Inspektor Emeritussignum}
Inspektor Emeritussignum utdelas av Medaljelelekommittén till D-sektionens avgående Inspektor vid installation av ny Inspektor för D-sektionen. Installationen sker under årligt ordenskapitel, eller vid annat av Medaljelelekommittén och D-sektionens sektionsmöte valt tillfälle. Inspektor Emeritussignum utdelas till en och samma person endast en gång.

\subsection{Hedersledamotsignum}
Hedersledamotsignum utdelas av Medaljelelekommittén till av D-sektionens sektionsmöte vald hedersledamot av D-sektionen vid installation av denne. Installationen sker under årligt ordenskapitel, eller vid annat av Medaljelelekommittén och D-sektionens sektionsmöte valt tillfälle. Hedersledamotsignum utdelas till en och samma person endast en gång.

\subsection{Medaljelelekommittémedaljen}
Medaljelelekommittémedaljen bäres av de som är medlemmar i Medaljelelekommittéen.

\subsection{DuCtig-medaljen}
DuCtig-medaljen erhålles som en del av studierådets utmärkelse DuCtig (Data und infocoms tentaprestationer i grundblocket). Medaljen finns i tre valörer; brons, silver och guld. Medaljen delas ut enligt studierådets kriterier.

\subsection{Styrelsemedaljen}
Styrelsemedaljen tilldelas de personer som innehar funktionärspost inom sektionsstyrelsen. En
styrelsemedalj delas ut vid påbörjad mandatperiod på posten. Styrelsemedaljen utdelas till en
och samma person endast en gång. Denna medalj tillhandages gratis.

\subsection{Insignier}
Insignier är små plaketter som fästes på erhållen funktionärsmedalj.

\begin{reglemlista}
	\item[Styrelseinsigniner]
	Styrelseinsigniner tilldelas de personer som innehaft funktionärspost inom sektionsstyrelsen. En styrelseinsigna delas ut för varje påbörjad mandatperiod på posten.
	\item[Gammal och äcklig]
	Utdelas till de som uppfyller kraven på Gammal och äcklig: nämligen att ha varit vald funktionär i tre hela år, samt innehaft ett mandat i sektionsstyrelsen eller att ha varit vald funktionär i fyra hela år
	\item[TLTH-Fullmäktige]
	Tilldelas den som valts in som ordinarie ledamot eller suppleant i Teknologkårens fullmäktige

\end{reglemlista}
\subsection{Utskottsmedaljer}
    
    Tilldelas de personer som varit vald funktionär inom samma utskott i tre hela år.  

\subsection{Rosa Pantern}
Rosa Pantern är ett diplom som utdelas av Studierådsordförande och Ordförande på årligt ordenskapitel, till person som de har funnit dels innehaft förtroendeuppdrag inom D-sektionen, dels på olika sätt verkat mycket föredömligt för D-sektionen, samt klarat överväldigande majoritet av sina studier på utsatt tid.

\subsection{Gyllene Pekpinnen}
Priset delas ut till en person, som har utvecklat pedagogiken på en eller flera kurser som läses av Datateknik, InfoCom, eller VRAR. Mottagaren kan vara doktorand men inte grundutbildningsstudent.

\end{document}



